
\documentclass[11pt,a4paper]{article}
\usepackage[margin=2.5cm]{geometry}
\usepackage[utf8]{inputenc}
\usepackage[T1]{fontenc}
\usepackage[french]{babel}
\usepackage{enumitem}
\usepackage{hyperref}


\begin{document}
	\pagestyle{empty}
	
	\noindent
	Friedly WOLI \\  % \hfill Airbus \\
	Toulouse, 31100  \\
	0767339290 \\
	friedlysedjro@gmail.com 
	
	\begin{flushright}
		
		Université Paris‑Saclay \\
		91190 Saint-Aubin, France
	\end{flushright}
	
	
	\bigskip
	
	
	
	
	
Objet : Candidature au Master 1 Méthodes Mathématiques pour la Mécanique \\
	
	Madame, Monsieur,\\
	\\
		Actuellement étudiant en mécanique énergétique à l’Université de Toulouse, je souhaite intégrer le Master 1 Méthodes Mathématiques pour la Mécanique de l’Université Paris-Saclay afin de renforcer mes compétences en modélisation mathématique et en analyse des phénomènes mécaniques complexes.
		
		Au cours de ma formation, j’ai acquis de solides bases en mécanique des milieux continus, en mécanique des fluides, en transferts thermiques ainsi qu’en mathématiques appliquées. J’ai développé un intérêt particulier pour les outils analytiques et numériques permettant de décrire rigoureusement les systèmes physiques. L’étude des équations aux dérivées partielles, de l’analyse fonctionnelle et des méthodes numériques m’a permis de mieux comprendre les fondements mathématiques des modèles utilisés en ingénierie.
		
		Le Master Méthodes Mathématiques pour la Mécanique représente pour moi une étape essentielle afin d’approfondir cette dimension théorique. La place centrale accordée aux fondements mathématiques de la mécanique, à la modélisation avancée et à l’analyse des modèles constitue un cadre idéal pour consolider ma rigueur scientifique et ma capacité d’abstraction. Je suis particulièrement motivé par l’exigence académique de cette formation et par son articulation entre théorie mathématique et applications mécaniques.
		
		Mon projet professionnel est de m’orienter vers la recherche en modélisation numérique du comportement des structures et matériaux, dans la perspective d’un doctorat. Je souhaite contribuer au développement de modèles prédictifs et de méthodes de calcul performantes adaptées aux problématiques industrielles, notamment dans les secteurs aéronautique et énergétique.
		
		Rigoureux, persévérant et doté d’un goût prononcé pour l’analyse, je suis prêt à m’investir pleinement dans cette formation exigeante. Je suis convaincu que le cadre académique stimulant de l’Université Paris-Saclay me permettra de développer les compétences scientifiques nécessaires pour atteindre mes objectifs.\\
		
		Je vous remercie par avance de l’attention que vous porterez à ma candidature et vous prie d’agréer, Madame, Monsieur, l’expression de mes salutations distinguées.
		
		
	
	
	\vspace{0.8cm}
	
	
	
	
	
	
	\begin{flushright}
		
		\textbf{Friedly WOLI}
	\end{flushright}
	
	
\end{document}




